\section{Data-Build Pipeline Instantiation}
\label{sec:pipeline}

The data-build pipeline instantiates the workflow abstraction as a sequence of artifact-producing stages. It begins with object concept expansion because the data engine needs explicit object identities, descriptions, and synthesis prompts before visual generation can be controlled. White-background object generation then produces isolated object images, which are easier to segment and reconstruct than images with scene clutter.

Foreground segmentation converts each generated object image into an RGBA asset suitable for image-to-3D reconstruction. This stage is aligned with promptable segmentation methods in the SAM family \cite{kirillov2023segmentanything,ravi2024sam2}, but the report treats segmentation as an implementation component rather than a new segmentation contribution. The reconstructed mesh gives the workflow an object that can be inserted into a scene, transformed, grounded, lit, and rendered from controlled viewpoints; image-to-3D asset-generation systems such as Hunyuan3D 2.1 provide relevant background for this stage \cite{teamhunyuan3d2025hunyuan3d21}. These stages are not independent scripts in the conceptual framework; they are artifact contracts that define what the controller can inspect and reuse.

Scene-aware rendering then places the reconstructed object into a fixed Blender environment \cite{blenderFoundationBlender}. This stage is where the original manual tuning pain point appears most strongly. The render must preserve the scene, use stable lighting, avoid floating or clipping, and keep the object at a plausible scale. VLM/CV review provides construction-time feedback, while controller actions adjust parameters, trigger rerendering, or mark samples for regeneration or filtering.

After a usable canonical source state is selected, the workflow exports target views and constructs train-ready source-target pairs. Each pair records the source image, target image, instruction, target rotation, object identifier, object name, and prompt version. The schema is intentionally close to the downstream image-editing interface: the source image and instruction form the input, and the target image provides the supervised edit target. This generic pair schema becomes concrete in the horizontal rotation case study, where the source is a canonical front view and each target is a controlled object-yaw view in the same scene. This design makes the dataset directly usable by the LoRA validation probe while keeping the construction trace inspectable.
